\chapter{Chat und Befehle}
\label{ch:players/chat}

Mods können mit dem Spielerchat arbeiten: Nachrichten senden,
Nachrichten abfangen und Chatbefehle\index{chatbefehl@Chatbefehl}
registrieren.

\section{Nachrichten senden}\label{nachrichten-senden}

\subsection{An alle Spieler}\label{an-alle-spieler}

Um eine Nachricht an alle Spieler im Spiel zu senden, rufst du
\texttt{chat\_\allowbreak{}send\_\allowbreak{}all} auf:

\begin{Shaded}
\begin{Highlighting}[]
\VariableTok{core}\OperatorTok{.}\NormalTok{chat\_send\_all}\OperatorTok{(}
    \StringTok{"Das ist eine Chatnachricht an alle Spieler"}
\OperatorTok{)}
\end{Highlighting}
\end{Shaded}

Im Spiel sieht das zum Beispiel so aus:

\begin{Shaded}
\begin{Highlighting}[]
\NormalTok{\textless{}player1\textgreater{} Schau dir diesen Eingang an}
\NormalTok{Das ist eine Chatnachricht an alle Spieler}
\NormalTok{\textless{}player2\textgreater{} Was ist damit?}
\end{Highlighting}
\end{Shaded}

Die Nachricht steht in einer eigenen Zeile, damit man sie vom Chat der
Spieler unterscheiden kann.

\subsection{An bestimmte Spieler}\label{an-bestimmte-spieler}

Um eine Nachricht an einen bestimmten Spieler zu senden, rufst du
\texttt{chat\_\allowbreak{}send\_\allowbreak{}player} auf:

\begin{Shaded}
\begin{Highlighting}[]
\VariableTok{core}\OperatorTok{.}\NormalTok{chat\_send\_player}\OperatorTok{(}
    \StringTok{"player1"}\OperatorTok{,}
    \StringTok{"Diese Nachricht ist für player1"}
\OperatorTok{)}
\end{Highlighting}
\end{Shaded}

Die Nachricht wird genauso angezeigt wie Nachrichten an alle Spieler,
ist aber nur für den genannten Spieler sichtbar, hier also für
\texttt{player1}.

\section{Chatbefehle}\label{chatbefehle}

Einen Chatbefehl, zum Beispiel \texttt{/\allowbreak{}foo}, registrierst
du mit \texttt{register\_\allowbreak{}chatcommand}:

\begin{Shaded}
\begin{Highlighting}[]
\VariableTok{core}\OperatorTok{.}\NormalTok{register\_chatcommand}\OperatorTok{(}\StringTok{"foo"}\OperatorTok{,} \OperatorTok{\{}
    \VariableTok{params} \OperatorTok{=} \StringTok{"\textless{}text\textgreater{}"}\OperatorTok{,}
    \VariableTok{description} \OperatorTok{=} \StringTok{"Wiederholt den Text"}\OperatorTok{,}
    \VariableTok{privs} \OperatorTok{=} \OperatorTok{\{}
        \VariableTok{interact} \OperatorTok{=} \KeywordTok{true}\OperatorTok{,}
    \OperatorTok{\},}
    \VariableTok{func} \OperatorTok{=} \KeywordTok{function}\OperatorTok{(}\VariableTok{name}\OperatorTok{,} \VariableTok{param}\OperatorTok{)}
        \ControlFlowTok{return} \KeywordTok{true}\OperatorTok{,} \StringTok{"Du hast "} \OperatorTok{..} \VariableTok{param} \OperatorTok{..} \StringTok{" gesagt!"}
    \KeywordTok{end}\OperatorTok{,}
\OperatorTok{\})}
\end{Highlighting}
\end{Shaded}

Hier ist \texttt{interact} als nötiges
Privileg\index{privileg@Privileg}~(Kapitel~\ref{ch:players/privileges})
angegeben; nur Spieler mit dem Privileg \texttt{interact} können den
Befehl also ausführen. \texttt{params} und \texttt{description}
erscheinen in der Hilfe, die \texttt{/\allowbreak{}help} anzeigt.

\texttt{param} ist ein String mit allem, was der Spieler nach dem Namen
des Befehls schreibt. Tippt jemand zum Beispiel
\texttt{/\allowbreak{}grantme one,two,three}, ist \texttt{param} gleich
\texttt{one,two,three}.

Chatbefehle können bis zu zwei Werte zurückgeben: zuerst einen
Wahrheitswert, der angibt, ob der Befehl erfolgreich war, und dann eine
Nachricht an den Benutzer.

\begin{hinweis}{Auch Offline-Spieler können Befehle ausführen}

Statt eines Spielerobjekts wird ein Spielername übergeben, weil Mods
Befehle im Namen von Spielern ausführen können, die nicht online sind.
Die IRC-Brücke erlaubt Spielern zum Beispiel, Befehle auszuführen, ohne
dem Spiel beizutreten.

Geh also nicht davon aus, dass der Spieler online ist. Das kannst du
prüfen, indem du schaust, ob
\texttt{core.\allowbreak{}get\_\allowbreak{}player\_\allowbreak{}by\_\allowbreak{}name}\index{core@\texttt{core}!get_player_by_name@\texttt{core.\allowbreak{}get\_\allowbreak{}player\_\allowbreak{}by\_\allowbreak{}name}}
einen Spieler zurückgibt.

\end{hinweis}

\subsection{Mehrere Argumente}\label{mehrere-argumente}

\texttt{param} enthält alle Argumente eines Chatbefehls in einem
einzigen String. Oft muss ein Befehl daraus mehrere Argumente
herauslösen. Dafür gibt es zwei Wege: die String-Funktion \texttt{split}
von Luanti oder Lua-Patterns.

\subsubsection{Mit string.split}\label{mit-string.split}

Mit \texttt{string.\allowbreak{}split(\textquotedbl{} \textquotedbl{})}
zerlegst du einen String in Wörter:

\begin{Shaded}
\begin{Highlighting}[]
\KeywordTok{local} \VariableTok{parts} \OperatorTok{=} \VariableTok{param}\OperatorTok{:}\NormalTok{split}\OperatorTok{(}\StringTok{" "}\OperatorTok{)}
\KeywordTok{local} \VariableTok{cmd} \OperatorTok{=} \VariableTok{parts}\OperatorTok{[}\DecValTok{1}\OperatorTok{]}

\ControlFlowTok{if} \VariableTok{cmd} \OperatorTok{==} \StringTok{"join"} \ControlFlowTok{then}
    \KeywordTok{local} \VariableTok{team\_name} \OperatorTok{=} \VariableTok{parts}\OperatorTok{[}\DecValTok{2}\OperatorTok{]}
    \VariableTok{team}\OperatorTok{.}\NormalTok{join}\OperatorTok{(}\VariableTok{name}\OperatorTok{,} \VariableTok{team\_name}\OperatorTok{)}
    \ControlFlowTok{return} \KeywordTok{true}\OperatorTok{,} \StringTok{"Team beigetreten!"}
\ControlFlowTok{elseif} \VariableTok{cmd} \OperatorTok{==} \StringTok{"max\_users"} \ControlFlowTok{then}
    \KeywordTok{local} \VariableTok{team\_name} \OperatorTok{=} \VariableTok{parts}\OperatorTok{[}\DecValTok{2}\OperatorTok{]}
    \KeywordTok{local} \VariableTok{max\_users} \OperatorTok{=} \FunctionTok{tonumber}\OperatorTok{(}\VariableTok{parts}\OperatorTok{[}\DecValTok{3}\OperatorTok{])}
    \ControlFlowTok{if} \VariableTok{team\_name} \KeywordTok{and} \VariableTok{max\_users} \ControlFlowTok{then}
        \ControlFlowTok{return} \KeywordTok{true}\OperatorTok{,}
            \StringTok{"Maximale Spielerzahl von "}
                \OperatorTok{..} \VariableTok{team\_name}
                \OperatorTok{..} \StringTok{" auf "}
                \OperatorTok{..} \VariableTok{max\_users}
                \OperatorTok{..} \StringTok{" gesetzt"}
    \ControlFlowTok{else}
        \ControlFlowTok{return} \KeywordTok{false}\OperatorTok{,}
            \StringTok{"Aufruf: /team max\_users \textless{}team\_name\textgreater{} \textless{}zahl\textgreater{}"}
    \ControlFlowTok{end}
\ControlFlowTok{else}
    \ControlFlowTok{return} \KeywordTok{false}\OperatorTok{,} \StringTok{"Befehl fehlt"}
\ControlFlowTok{end}
\end{Highlighting}
\end{Shaded}

\subsubsection{Mit Lua-Patterns}\label{mit-lua-patterns}

Mit Lua-Patterns\footnote{\url{https://www.lua.org/pil/20.2.html}}
(englisch) ziehst du anhand von Regeln Teile aus einem Text heraus. Sie
eignen sich am besten, wenn Argumente Leerzeichen enthalten können oder
du genauer steuern willst, wie Parameter erfasst werden.

\begin{Shaded}
\begin{Highlighting}[]
\KeywordTok{local} \VariableTok{to}\OperatorTok{,} \VariableTok{msg} \OperatorTok{=} \VariableTok{param}\OperatorTok{:}\FunctionTok{match}\OperatorTok{(}\StringTok{"\^{}([\%a\%d\_{-}]+) (.+)$"}\OperatorTok{)}
\end{Highlighting}
\end{Shaded}

Dieser Code setzt \texttt{/\allowbreak{}msg <an> <nachricht>} um. Gehen
wir ihn von links nach rechts durch:

\begin{itemize}
\tightlist
\item
  \texttt{\textasciicircum{}} passt auf den Anfang des Strings.
\item
  \texttt{()} ist eine Gruppe\index{gruppe@Gruppe} -- was darin passt,
  gibt \texttt{string.\allowbreak{}match} zurück.
\item
  \texttt{[]} akzeptiert die Zeichen in dieser Liste.
\item
  \texttt{\%a} akzeptiert jeden Buchstaben, \texttt{\%d} jede Ziffer.
\item
  \texttt{[\%a\%d\_\allowbreak{}-]} akzeptiert also jeden Buchstaben,
  jede Ziffer, \texttt{\_\allowbreak{}} und \texttt{-}.
\item
  \texttt{+} passt, wenn das Vorherige ein- oder mehrmals vorkommt.
\item
  \texttt{.\allowbreak{}} passt hier auf jedes beliebige Zeichen.
\item
  \texttt{\$} passt auf das Ende des Strings.
\end{itemize}

Kurz gesagt: Das Pattern erfasst den Namen (ein Wort nur aus Buchstaben,
Ziffern, \texttt{-} und \texttt{\_\allowbreak{}}), dann ein Leerzeichen,
dann die Nachricht (ein oder mehr beliebige Zeichen). Name und Nachricht
werden zurückgegeben, weil sie in Klammern stehen.

So setzen die meisten Mods komplexe Chatbefehle um. Ausführlichere
Anleitungen zu Lua-Patterns sind das Tutorial auf
lua-users.org\footnote{\url{http://lua-users.org/wiki/PatternsTutorial}}
(englisch) und die Dokumentation in PIL\footnote{\url{https://www.lua.org/pil/20.2.html}}
(englisch).

\section{Nachrichten abfangen}\label{nachrichten-abfangen}

Um Nachrichten abzufangen, verwendest du
\texttt{register\_\allowbreak{}on\_\allowbreak{}chat\_\allowbreak{}message}:

\begin{Shaded}
\begin{Highlighting}[]
\VariableTok{core}\OperatorTok{.}\NormalTok{register\_on\_chat\_message}\OperatorTok{(}\KeywordTok{function}\OperatorTok{(}\VariableTok{name}\OperatorTok{,} \VariableTok{message}\OperatorTok{)}
    \FunctionTok{print}\OperatorTok{(}\VariableTok{name} \OperatorTok{..} \StringTok{" sagt "} \OperatorTok{..} \VariableTok{message}\OperatorTok{)}
    \ControlFlowTok{return} \KeywordTok{false}
\KeywordTok{end}\OperatorTok{)}
\end{Highlighting}
\end{Shaded}

Mit \texttt{return false} lässt du zu, dass die Nachricht wie üblich
gesendet wird. Du könntest die Zeile \texttt{return false} auch
weglassen, ohne dass sich etwas ändert, denn dann wird implizit
\texttt{nil} zurückgegeben, und das wird wie \texttt{false} behandelt.

\begin{hinweis}{Privilegien und Chatbefehle}

Ein Spieler braucht das Privileg \texttt{shout} nicht, um diesen
Callback\index{callback@Callback} auszulösen. Chatbefehle sind nämlich
in Lua umgesetzt und einfach Chatnachrichten, die mit
\texttt{/\allowbreak{}} beginnen.

\end{hinweis}

Berücksichtige also, dass es sich um einen Chatbefehl handeln kann oder
dass dem Spieler \texttt{shout} fehlt.

\begin{Shaded}
\begin{Highlighting}[]
\VariableTok{core}\OperatorTok{.}\NormalTok{register\_on\_chat\_message}\OperatorTok{(}\KeywordTok{function}\OperatorTok{(}\VariableTok{name}\OperatorTok{,} \VariableTok{message}\OperatorTok{)}
    \ControlFlowTok{if} \VariableTok{message}\OperatorTok{:}\FunctionTok{sub}\OperatorTok{(}\DecValTok{1}\OperatorTok{,} \DecValTok{1}\OperatorTok{)} \OperatorTok{==} \StringTok{"/"} \ControlFlowTok{then}
        \FunctionTok{print}\OperatorTok{(}\VariableTok{name} \OperatorTok{..} \StringTok{" hat einen Chatbefehl ausgeführt"}\OperatorTok{)}
    \ControlFlowTok{elseif}
        \VariableTok{core}\OperatorTok{.}\NormalTok{check\_player\_privs}\OperatorTok{(}\VariableTok{name}\OperatorTok{,} \OperatorTok{\{} \VariableTok{shout} \OperatorTok{=} \KeywordTok{true} \OperatorTok{\})}
    \ControlFlowTok{then}
        \FunctionTok{print}\OperatorTok{(}\VariableTok{name} \OperatorTok{..} \StringTok{" sagt "} \OperatorTok{..} \VariableTok{message}\OperatorTok{)}
    \ControlFlowTok{else}
        \FunctionTok{print}\OperatorTok{(}
            \VariableTok{name}
                \OperatorTok{..} \StringTok{" wollte "}
                \OperatorTok{..} \VariableTok{message}
                \OperatorTok{..} \StringTok{" sagen, hat aber kein shout"}
        \OperatorTok{)}
    \ControlFlowTok{end}

    \ControlFlowTok{return} \KeywordTok{false}
\KeywordTok{end}\OperatorTok{)}
\end{Highlighting}
\end{Shaded}
